% vision ?
\section{ArmoniK ambitions are to provide a high-performance computing platform for the most demanding applications.}


\begin{frame}{What is ArmoniK?}
	\begin{alertblock}{ArmoniK}
		A task orchestrator.
	\end{alertblock}
\end{frame}

\begin{frame}{ArmoniK's envisioned capabilities}
	\begin{itemize}
		\item \textbf{Production ready}, with support for security, scalability, and reliability.
		\item \textbf{Easy to use and deploy}, with a user-friendly interface and comprehensive documentation.
		\item \textbf{Fault tolerant}, with support for automatic recovery and rescheduling of failed tasks.
		\item \textbf{Observable}, with support for monitoring and logging of task execution.
		\item \textbf{Malleable}, with support for dynamic resource allocation.
		\item \textbf{Portable}, with support for multiple platforms and cloud providers.
		\item \textbf{Data aware}, with support for data locality and efficient data management.
		\item \textbf{Performant}, with support for high-throughput and high-volume task execution.
		\item \textbf{Smart scheduling}, with priority, preemption, and heterogeneous resource placement.
		\item \textbf{Open source}, with a vibrant community of users and contributors.
	\end{itemize}
\end{frame}


\begin{frame}{Production Ready}
	\begin{itemize}
		\item \textbf{Security by Architecture:} Strict separation of Control, Data, and Compute Planes
		\item \textbf{Data Protection:} Encryption at rest and in transit, secure secret injection for workers
		\item \textbf{High Availability:} No single point of failure, automated failover and recovery
		\item \textbf{Horizontal Scalability:} Stateless components scale out under high concurrency
		\item \textbf{Robust Error Handling:} Idempotent operations with well-defined state transitions
		\item \textbf{Multi-Tenancy:} Tenant isolation, per-tenant quotas, and noisy-neighbor prevention
		\item \textbf{Enterprise Governance:} Identity integration, role separation, and auditability
	\end{itemize}
\end{frame}

\begin{frame}{Easy to Use \& Deploy}
	\begin{itemize}
		\item \textbf{Flexible Deployment:} From fully managed to self-operated
		\item \textbf{Low-Friction Onboarding:} Automated provisioning and standardized configurations
		\item \textbf{Declarative Configuration:} Repeatable, auditable, drift-free deployments
		\item \textbf{Developer-Centric SDKs:} Multi-language APIs that abstract orchestration complexity
		\item \textbf{Local Dev \& Test:} Run and debug tasks locally before deploying to the cluster
		\item \textbf{Worker Versioning:} Canary rollouts, rollback, and side-by-side version coexistence
		\item \textbf{Reduced Operational Burden:} Clear responsibilities and comprehensive documentation
	\end{itemize}
\end{frame}

\begin{frame}{Fault Tolerant}
	\begin{itemize}
		\item \textbf{Automatic Recovery:} Failed tasks are retried or rescheduled transparently
		\item \textbf{Checkpointing:} Persist task execution state for efficient resumption after failures
		\item \textbf{Graceful Degradation:} Services restart without impacting running workloads
		\item \textbf{Failure Isolation:} Contained within regions, partitions, and execution domains
		\item \textbf{Idempotent Operations:} Safe retries with consistent outcomes
	\end{itemize}
\end{frame}

\begin{frame}{Observable}
	\begin{itemize}
		\item \textbf{End-to-End Visibility:} Telemetry across scheduling, execution, and system health
		\item \textbf{Functional Observability:} Task graphs, critical paths, and execution structure
		\item \textbf{Unified View:} Correlated data across control, compute, and data domains
		\item \textbf{Business Metrics:} Domain-specific tags bridging execution and business outcomes
		\item \textbf{Intelligent Insight:} Automated anomaly detection and diagnostics
		\item \textbf{Custom Instrumentation:} User-defined metrics and tracing for domain-specific observability, also available to tasks during execution to adapt behavior based on runtime conditions
		\item \textbf{Metering \& Cost Attribution:} Per-tenant usage tracking for billing and chargeback
	\end{itemize}
\end{frame}

\begin{frame}{Malleable}
	\begin{itemize}
		\item \textbf{Elastic Resources:} Scale execution capacity up or down with workload demand
		\item \textbf{Runtime Reconfiguration:} Resize execution domains without restarts
		\item \textbf{Fixed Infrastructure:} Adjust active resource usage to avoid idle capacity
		\item \textbf{Policy-Driven:} Balance performance, cost, and fairness objectives
		\item \textbf{Predictive Adaptation:} Proactive scaling based on workload patterns
        \item \textbf{Resource Sharing:} Seamless workload distribution across environments based on resource availability and policy constraints
	\end{itemize}
\end{frame}

\begin{frame}{Portable}
	\begin{itemize}
		\item \textbf{Multi-Environment:} Cloud, on-premises, bare-metal, and supercomputing
		\item \textbf{Infrastructure-Agnostic:} Resources treated as a logical pool
		\item \textbf{Hybrid \& Multi-Cloud:} Distribute workloads across heterogeneous environments
		\item \textbf{Standards-Oriented:} Open standards to reduce ecosystem lock-in
		\item \textbf{Data Sovereignty:} Compliance-aware placement respecting residency requirements
	\end{itemize}
\end{frame}

\begin{frame}{Data Aware}
	\begin{itemize}
		\item \textbf{Data-Driven Scheduling:} Task placement decisions account for data locality, dependencies, and reuse to minimize data movement
		\item \textbf{Proactive Placement:} Replicate or stage data ahead of execution based on the task graph structure
		\item \textbf{Flexible Dependencies:} Mandatory, optional, and on-demand data classes enabling fine-grained scheduling trade-offs
		\item \textbf{Controlled Lifecycle:} Explicit ownership, locality, and retention management
		\item \textbf{Cross-Cluster Migration:} Move data and reschedule tasks across environments for performance, cost, and compliance
	\end{itemize}
\end{frame}

\begin{frame}{Performant}
	\begin{itemize}
		\item \textbf{Low-Latency Orchestration:} Minimized coordination overhead on the critical path
		\item \textbf{Pipeline Parallelism:} Overlap data transfers and computation
		\item \textbf{High Concurrency:} Sustained throughput with large worker counts
		\item \textbf{Efficient Utilization:} Autoscaling to maximize useful work, minimize waste
		\item \textbf{Decoupled Architecture:} Loose coupling for throughput and fault isolation
	\end{itemize}
\end{frame}

\begin{frame}{Smart Scheduling}
	\begin{itemize}
		\item \textbf{Data-Aware Placement:} Co-locate tasks with their data to reduce transfer overhead
		\item \textbf{Priority \& Preemption:} Urgent tasks preempt batch work with configurable policies
		\item \textbf{Heterogeneous Compute:} GPU, accelerator, and NUMA-aware resource matching
		\item \textbf{Pluggable Policies:} Fair-share, backfill, gang scheduling, deadline-aware
		\item \textbf{DAG Optimizations:} Critical-path scheduling, task fusion, and speculative execution
		\item \textbf{Cost-Aware Decisions:} Task duration and resource prediction inform placement and scaling
	\end{itemize}
\end{frame}

\begin{frame}{Open Source}
	\begin{itemize}
		\item \textbf{Transparency:} Inspect, audit, and understand platform behavior
		\item \textbf{Extensibility:} Adapt to domain-specific needs and custom integrations
		\item \textbf{Open Ecosystems:} Standards-based portability and long-term sustainability
		\item \textbf{Community-Driven:} Shared innovation, faster iteration, no vendor lock-in
	\end{itemize}
\end{frame}